% Отчёт по практикуму. Сборка: pdflatex report.tex (дважды — для оглавления)
\documentclass[12pt,a4paper]{article}

\usepackage[utf8]{inputenc}
\usepackage[T2A]{fontenc}
\usepackage[russian]{babel}
\usepackage[left=25mm,right=15mm,top=20mm,bottom=20mm]{geometry}
\usepackage{indentfirst}
\usepackage{fancyvrb}
\usepackage{booktabs}
\usepackage{float}
\usepackage{graphicx}
\usepackage{tikz}
\usetikzlibrary{arrows.meta,positioning}
\usepackage[hidelinks]{hyperref}

\fvset{fontsize=\small,frame=single,framesep=2mm}

% Обозначения для диаграмм
\tikzset{
  cls/.style={draw,rectangle,minimum width=2.9cm,minimum height=0.8cm,
              font=\small\ttfamily,align=center,fill=white},
  obj/.style={draw,rounded corners=7pt,minimum width=2.9cm,minimum height=0.8cm,
              font=\small\ttfamily,align=center,fill=white},
  inh/.style={-{Triangle[open,length=8pt,width=8pt]}},
  agg/.style={{Diamond[open,length=10pt,width=6pt]}-},
  comp/.style={{Diamond[length=10pt,width=6pt]}-},
  use/.style={dashed,-{Stealth[length=6pt]}},
  msg/.style={-{Stealth[length=6pt]}},
  lbl/.style={font=\scriptsize,fill=white,inner sep=1pt},
}

\begin{document}

% ---------- Титульный лист ----------
\begin{titlepage}
\centering
\vfill
{\large Отчёт по практикуму\\ по объектно-ориентированному программированию}\\[1cm]
{\LARGE\bfseries Система автоматизации функций секретаря}
\vfill
\begin{flushright}
Выполнили: студенты группы \underline{\hspace{2cm}}\\
Волков Георгий\\
Усков Никита
\end{flushright}
\vfill
Москва, 2026
\end{titlepage}

\tableofcontents
\newpage

% =====================================================================
\section{Уточнённая постановка задачи}

Фирма состоит из $K$ отделов ($5 \le K \le 9$). В каждом отделе есть руководитель
и несколько сотрудников, а также своё помещение. Кроме того, у фирмы есть общий
конференц-зал. Работа отделов согласуется с помощью \textbf{общего календаря}
событий. Виды событий:
\begin{itemize}
  \item ежедневная планёрка отдела (в начале рабочего дня, в помещении отдела);
  \item совещание руководителей отделов (по понедельникам, в конференц-зале);
  \item тематический семинар для сотрудников разных отделов (по средам);
  \item рабочее совещание группы сотрудников;
  \item командировка сотрудника (одна или несколько дней, вне фирмы).
\end{itemize}

У каждого события есть название, время начала, длительность, место проведения,
список участников и степень важности (низкая, средняя, высокая).

Система автоматизирует следующие функции секретаря:
\begin{enumerate}
  \item \textbf{Актуализация календаря}: прошедшие события удаляются автоматически;
        пользователь может добавлять новые события, изменять и удалять существующие.
  \item \textbf{Просмотр календаря} на выбранный день для всей фирмы,
        для отдела или для конкретного сотрудника.
  \item \textbf{Напоминания} участникам о предстоящих событиях (моделируются как
        отправленные письма). Начало и частота напоминаний зависят от важности
        события; по умолчанию: низкая --- один раз за 1~день, средняя --- каждый
        день за 2~дня, высокая --- каждый день за 4~дня.
  \item \textbf{Контроль конфликтов} при добавлении и изменении событий.
\end{enumerate}

\subsection*{Конфликты и их разрешение}
Выявляются три вида конфликтов:
\begin{itemize}
  \item сотрудник в это время уже участвует в другом событии;
  \item помещение в это время уже занято;
  \item участников больше, чем мест в помещении.
\end{itemize}
Планёрку разрешается пропустить: если у сотрудника в это время другое событие
(например, командировка), он просто исключается из участников планёрки,
и конфликтом это не считается.

При обнаружении конфликта система подбирает варианты его разрешения:
\begin{itemize}
  \item перенести новое событие на ближайшее свободное время;
  \item провести новое событие в другом помещении;
  \item не приглашать занятых сотрудников;
  \item перенести уже назначенное событие.
\end{itemize}
Пользователь выбирает один из вариантов, исправляет событие вручную или
отклоняет заявку. В системе две категории пользователей: \textbf{администратор}
может изменять и удалять уже назначенные события (в том числе выбирать последний
вариант), \textbf{обычный пользователь} --- нет.

\subsection*{Моделирование}
Задаётся начальное состояние календаря (регулярные события на весь период и две
командировки), после чего время изменяется пошагово. Период моделирования ---
$N$ дней ($7 \le N \le 30$), шаг --- 30 минут или 1 час. Рабочий день по умолчанию
с 9:00 до 18:00; ночи и выходные пропускаются.

На каждом шаге:
\begin{enumerate}
  \item из календаря удаляются завершившиеся события;
  \item рассылаются напоминания;
  \item с некоторой вероятностью поступает новая заявка на событие, которая
        проверяется на конфликты.
\end{enumerate}

Случайные заявки моделируются статистически. Если в среднем за рабочий день
поступает $\lambda$ заявок, а длина рабочего дня $T$, то на шаге $\Delta t$ заявка
поступает с вероятностью $p = \lambda \Delta t / T$. Остальные характеристики
заявки выбираются по равномерному закону:
\begin{itemize}
  \item день проведения --- от следующего дня до пяти дней вперёд;
  \item время начала --- любое получасовое время рабочего дня;
  \item вид события --- совещание 45\,\%, семинар 25\,\%, командировка 15\,\%,
        внеплановое совещание руководителей 15\,\%;
  \item важность --- низкая 30\,\%, средняя 50\,\%, высокая 20\,\%;
  \item число участников --- в заданных для каждого вида пределах.
\end{itemize}

Конфликты случайных заявок разрешаются либо пользователем (моделирование
приостанавливается), либо автоматически --- выбирается первый допустимый
для текущей категории пользователя вариант.

\textbf{Параметры моделирования} (задаются перед каждым экспериментом):
\begin{itemize}
  \item $K$ и число сотрудников в отделе;
  \item $N$ и шаг моделирования;
  \item среднее число заявок в день $\lambda$ и начальное значение генератора
        случайных чисел;
  \item режим разрешения конфликтов и категория пользователя.
\end{itemize}
Дата начала --- текущая дата, рабочие часы и правила напоминаний фиксированы
(см. выше).
По окончании моделирования выводится итоговый отчёт:
\begin{itemize}
  \item список состоявшихся событий;
  \item число заявок, конфликтов, отклонённых заявок и напоминаний;
  \item загрузка конференц-зала.
\end{itemize}

% =====================================================================
\section{Диаграмма классов}

Ниже приведены три диаграммы. На рис.~\ref{fig:hier} показаны иерархии
наследования, на рис.~\ref{fig:domain} --- связи классов предметной области,
на рис.~\ref{fig:struct} --- общая структура моделирующей части системы.

Обозначения на диаграммах:
\begin{itemize}
  \item стрелка с полым треугольником --- наследование;
  \item полый ромб --- агрегация;
  \item закрашенный ромб --- композиция;
  \item пунктирная стрелка --- использование.
\end{itemize}

\begin{figure}[H]
\centering
\begin{tikzpicture}[node distance=1.2cm]
  % События
  \node[cls] (event) at (0,0) {Event\\\textit{(абстрактный)}};
  \node[cls,minimum width=2.6cm,font=\footnotesize\ttfamily] (dm)  at (-6.4,-2.2) {DailyMeeting};
  \node[cls,minimum width=2.6cm,font=\footnotesize\ttfamily] (mm)  at (-3.2,-2.2) {ManagersMeeting};
  \node[cls,minimum width=2.6cm,font=\footnotesize\ttfamily] (sem) at (0,-2.2)    {Seminar};
  \node[cls,minimum width=2.6cm,font=\footnotesize\ttfamily] (wm)  at (3.2,-2.2)  {WorkMeeting};
  \node[cls,minimum width=2.6cm,font=\footnotesize\ttfamily] (bt)  at (6.4,-2.2)  {BusinessTrip};
  \foreach \c in {dm,mm,sem,wm,bt} \draw[inh] (\c.north) -- (event.south);

  % Люди
  \node[cls] (person) at (-4,-4.4) {Person\\\textit{(абстрактный)}};
  \node[cls] (emp)    at (-4,-6.2) {Employee};
  \node[cls] (head)   at (-4,-8.0) {DepartmentHead};
  \draw[inh] (emp.north) -- (person.south);
  \draw[inh] (head.north) -- (emp.south);

  % Помещения
  \node[cls] (room)  at (3,-4.4) {Room\\\textit{(абстрактный)}};
  \node[cls] (droom) at (1.2,-6.6) {DepartmentRoom};
  \node[cls] (hall)  at (4.8,-6.6) {ConferenceHall};
  \draw[inh] (droom.north) -- (room.south);
  \draw[inh] (hall.north) -- (room.south);
\end{tikzpicture}
\caption{Иерархии наследования: события, люди, помещения}
\label{fig:hier}
\end{figure}

\begin{figure}[H]
\centering
\begin{tikzpicture}
  \node[cls] (company) at (0,0)     {Company};
  \node[cls] (dept)    at (0,-2.2)  {Department};
  \node[cls] (emp)     at (0,-4.4)  {Employee};
  \node[cls] (room)    at (-5,-4.4) {Room};
  \node[cls] (cal)     at (5,0)     {Calendar};
  \node[cls] (event)   at (5,-2.2)  {Event};

  \draw[comp] (company.south) -- node[lbl,pos=0.75,right] {1..K} (dept.north);
  \draw[comp] (dept.south) -- node[lbl,pos=0.75,right] {1..*} (emp.north);
  \draw[comp] (company.west) -| node[lbl,pos=0.9,left] {1..*} (room.north);
  \draw[comp] (cal.south) -- node[lbl,pos=0.75,right] {0..*} (event.north);
  \draw[agg]  (event.south west) -- node[lbl,pos=0.8,above] {участники, M:N} (emp.north east);
  \draw[use]  (event.south) -- ++(0,-3.0) -| node[lbl,pos=0.25,above] {место проведения} (room.south);
\end{tikzpicture}
\caption{Классы предметной области и связи между ними}
\label{fig:domain}
\end{figure}

\begin{figure}[H]
\centering
\begin{tikzpicture}
  \node[cls] (sim)  at (0,0)       {Simulation};
  \node[cls] (comp) at (-5.6,0)    {Company};
  \node[cls] (cal)  at (5.6,0)     {Calendar};
  \node[cls] (gen)  at (-2.8,2.8)  {EventGenerator};
  \node[cls] (not)  at (2.8,2.8)   {Notifier};
  \node[cls] (res)  at (-2.8,-2.8) {ConflictResolver};
  \node[cls] (chk)  at (2.8,-2.8)  {ConflictChecker};

  % Simulation содержит все части
  \draw[comp] (sim.west) -- (comp.east);
  \draw[comp] (sim.east) -- (cal.west);
  \draw[comp] (sim.north west) -- (gen.south);
  \draw[comp] (sim.north east) -- (not.south);
  \draw[comp] (sim.south west) -- (res.north);
  \draw[comp] (sim.south east) -- (chk.north);
  % ConflictResolver содержит свой ConflictChecker
  \draw[comp] (res.east) -- (chk.west);

  % Использование
  \draw[use] (gen.west) -| (comp.north);
  \draw[use] (gen.north) -- ++(0,0.6) -| (cal.north);
  \draw[use] (not.east) -| (cal.north);
  \draw[use] (chk.east) -| (cal.south);
  \draw[use] (res.west) -| (comp.south);
\end{tikzpicture}
\caption{Структура моделирующей части системы}
\label{fig:struct}
\end{figure}

Классы интерфейса показаны на рис.~\ref{fig:gui}. Они наследуются от
классов библиотеки Qt и работают с моделью только через объект
\texttt{Simulation}.

\begin{figure}[H]
\centering
\begin{tikzpicture}
  \node[cls] (qmw) at (-5,2.4)  {QMainWindow\\\textit{(Qt)}};
  \node[cls] (qdlg) at (2.6,2.4) {QDialog\\\textit{(Qt)}};
  \node[cls] (mw)  at (-5,0)    {MainWindow};
  \node[cls,minimum width=2.6cm,font=\footnotesize\ttfamily] (pd) at (-1.6,0) {ParamsDialog};
  \node[cls,minimum width=2.6cm,font=\footnotesize\ttfamily] (ed) at (1.2,0)  {EventDialog};
  \node[cls,minimum width=2.6cm,font=\footnotesize\ttfamily] (cd) at (4.0,0)  {ConflictDialog};
  \node[cls,minimum width=2.6cm,font=\footnotesize\ttfamily] (rd) at (6.8,0)  {ReportDialog};
  \node[cls] (sim) at (1.2,-2.6) {Simulation};

  \draw[inh] (mw.north) -- (qmw.south);
  \foreach \c in {pd,ed,cd,rd} \draw[inh] (\c.north) -- (qdlg.south);
  \draw[comp] (mw.south) |- (sim.west);
  \draw[use] (ed.south) -- (sim.north);
  \draw[use] (cd.south) -- (sim.north east);
  \draw[use] (rd.south) |- (sim.east);
  \draw[use] (mw.north east) -- ++(0.6,0.5) -| (pd.north);
\end{tikzpicture}
\caption{Классы пользовательского интерфейса}
\label{fig:gui}
\end{figure}

Вспомогательные структуры (на диаграммах не показаны):
\begin{itemize}
  \item \texttt{Conflict} --- описание найденного конфликта;
  \item \texttt{FixOption} --- вариант его разрешения;
  \item \texttt{PendingRequest} --- заявка, ожидающая решения пользователя;
  \item \texttt{SimulationParams} --- параметры моделирования;
  \item \texttt{Statistics} --- итоговая статистика;
  \item \texttt{Clock} и \texttt{WorkSchedule} --- перевод модельного времени
        в даты и рабочий график.
\end{itemize}

% =====================================================================
\section{Текстовые спецификации основных классов}

Модельное время хранится как целое число минут от 00:00 первого дня периода
(тип \texttt{Minutes}).

\subsection*{Event --- календарное событие}
\begin{Verbatim}
// Абстрактный базовый класс всех событий
class Event {
public:
    Event(string title, Minutes start, Minutes duration,
          Room* room, Importance importance);
    virtual EventKind kind() const = 0;              // вид события
    virtual unique_ptr<Event> clone() const = 0;      // копия события
    virtual bool allowsAbsence() const;               // можно ли пропустить
    virtual const Department* ownerDepartment() const;// отдел события
    // чтение: id(), title(), start(), end(), duration(), room(),
    //         importance(), participants()
    // изменение: setStart(), setDuration(), setRoom(), setImportance(),
    //            setParticipants(), removeParticipant()
    bool overlaps(const Event& other) const;  // пересечение по времени
    bool hasParticipant(const Employee* e) const;
    bool involves(const Department* d) const;
    void markReminded(Minutes t);              // учёт напоминаний
};
// Наследники: DailyMeeting (allowsAbsence() = true), ManagersMeeting,
//             Seminar, WorkMeeting, BusinessTrip (без помещения)
\end{Verbatim}

\subsection*{Calendar --- общий календарь фирмы}
\begin{Verbatim}
class Calendar {
public:
    int add(unique_ptr<Event> event);          // добавить, вернуть id
    unique_ptr<Event> remove(int id);          // убрать событие
    Event* find(int id) const;
    // убрать и вернуть события, закончившиеся к моменту now
    vector<unique_ptr<Event>> removeFinished(Minutes now);
    // события в промежутке [from, to) для фирмы / отдела / сотрудника
    vector<Event*> query(Minutes from, Minutes to,
                         const EventFilter& filter) const;
private:
    void markAbsences(Event& added);  // исключение из пропускаемых планёрок
    vector<unique_ptr<Event>> events_;  // календарь владеет событиями
};
\end{Verbatim}

\subsection*{Company --- фирма}
\begin{Verbatim}
class Company {
public:
    // создание K отделов со случайными сотрудниками
    static unique_ptr<Company> generate(int departments,
                                        int staffPerDepartment, mt19937& rng);
    const vector<unique_ptr<Department>>& departments() const;
    const vector<unique_ptr<Room>>& rooms() const;
    ConferenceHall* conferenceHall() const;
    vector<Employee*> allEmployees() const;
    vector<Employee*> heads() const;          // руководители отделов
};
\end{Verbatim}

\subsection*{ConflictChecker и ConflictResolver --- контроль конфликтов}
\begin{Verbatim}
class ConflictChecker {
public:
    // конфликты события с календарём (skipId - не учитывать это событие)
    vector<Conflict> check(const Event& event, int skipId = 0) const;
};

class ConflictResolver {
public:
    // варианты разрешения найденных конфликтов
    vector<FixOption> suggest(const Event& newEvent,
                              const vector<Conflict>& conflicts,
                              Minutes now, int skipId = 0) const;
    void apply(const FixOption& option, Event& newEvent);  // применить
    // номер первого допустимого варианта для авто-режима (или -1)
    int chooseAuto(const vector<FixOption>& options, bool isAdmin) const;
    // ближайшее свободное время не раньше from (или -1)
    Minutes findFreeTime(const Event& event, Minutes from, int skipId,
                         const Event* avoid = nullptr) const;
};
\end{Verbatim}

\subsection*{Notifier --- рассылка напоминаний}
\begin{Verbatim}
class Notifier {
public:
    explicit Notifier(const ReminderRule rules[3]);  // по степени важности
    // разослать напоминания, которые пора отправить; вернуть письма
    vector<Message> process(Calendar& calendar, Minutes now);
};
\end{Verbatim}

\subsection*{EventGenerator --- генерация событий}
\begin{Verbatim}
class EventGenerator {
public:
    // начальный календарь: планёрки, совещания, семинары, командировки
    void fillInitial(Calendar& calendar, const ConflictChecker& checker);
    bool requestArrives(double perDay, Minutes step);  // пришла ли заявка
    unique_ptr<Event> randomRequest(Minutes now);      // случайная заявка
};
\end{Verbatim}

\subsection*{Simulation --- моделирование}
\begin{Verbatim}
class Simulation {
public:
    explicit Simulation(const SimulationParams& params);
    void step();                 // один шаг моделирования
    bool finished() const;
    // заявки пользователя и разрешение конфликтов
    void submit(unique_ptr<Event> event, int replacesId = 0);
    bool hasPending() const;
    const PendingRequest& pending() const;
    void choosePending(int optionIndex);
    void replacePending(unique_ptr<Event> edited);
    void rejectPending();
    bool cancelEvent(int id);    // только администратор
    void setRole(UserRole role);
    // доступ: company(), calendar(), log(), messages(), history(), stats()
};
\end{Verbatim}

% =====================================================================
\section{Диаграмма объектов}

На рис.~\ref{fig:objects} показаны основные объекты системы во время
моделирования. Стрелки помечены вызываемыми операциями.

\begin{figure}[H]
\centering
\begin{tikzpicture}
  \node[obj] (mw)   at (0,4.8)     {mainWindow\\: MainWindow};
  \node[obj] (sim)  at (0,0)       {sim\\: Simulation};
  \node[obj] (comp) at (-5.6,0)    {company\\: Company};
  \node[obj] (cal)  at (5.6,0)     {calendar\\: Calendar};
  \node[obj] (gen)  at (-2.8,2.8)  {generator\\: EventGenerator};
  \node[obj] (not)  at (2.8,2.8)   {notifier\\: Notifier};
  \node[obj] (res)  at (-2.8,-2.8) {resolver\\: ConflictResolver};
  \node[obj] (chk)  at (2.8,-2.8)  {checker\\: ConflictChecker};

  \node[obj] (d1)  at (-5.6,-4.6) {sales\\: Department};
  \node[obj] (e1)  at (1.6,-4.6)  {petrov\\: Employee};
  \node[obj] (ev1) at (5.6,-4.6)  {seminar\\: Seminar};

  \draw[msg] (mw) -- node[lbl,right] {step(), submit()} (sim);
  \draw[msg] (sim) -- node[lbl,sloped,above] {randomRequest()} (gen);
  \draw[msg] (sim) -- node[lbl,sloped,above] {process()} (not);
  \draw[msg] (sim) -- node[lbl,sloped,below] {suggest(), apply()} (res);
  \draw[msg] (sim) -- node[lbl,sloped,below] {check()} (chk);
  \draw[msg] (sim) -- node[lbl,above] {add()} (cal);
  \draw[msg] (gen) -- node[lbl,sloped,above] {allEmployees()} (comp);
  \draw[msg] (not) -- node[lbl,sloped,above] {events()} (cal);
  \draw[msg] (chk) -- node[lbl,sloped,below] {events()} (cal);
  \draw[msg] (res) -- node[lbl,above] {check()} (chk);
  \draw[msg] (comp) -- node[lbl,right] {staff()} (d1);
  \draw[msg] (cal) -- node[lbl,right,align=left] {overlaps(),\\hasParticipant()} (ev1);
  \draw[msg] (ev1) -- node[lbl,above] {name()} (e1);
\end{tikzpicture}
\caption{Диаграмма объектов}
\label{fig:objects}
\end{figure}

Объект \texttt{mainWindow} по таймеру вызывает \texttt{step()}, а при
добавлении события --- \texttt{submit()}. Диалоги (окно события, окно
конфликта, отчёт) создаются главным окном и обращаются к тому же объекту
\texttt{sim}.

% =====================================================================
\section{Инструментальные средства}

\begin{itemize}
  \item Язык разработки --- C++ (стандарт C++20).
  \item Компилятор --- GCC 13.3.
  \item Система сборки --- CMake 3.28.
  \item Библиотека интерфейса --- Qt 5.15 (модуль Widgets).
  \item Операционная система --- Ubuntu 24.04.
  \item Среда разработки --- \underline{\hspace{4cm}}.
\end{itemize}

% =====================================================================
\section{Файловая структура системы}

Моделирующая часть (каталог \texttt{src/model}) не зависит от Qt и собирается
как отдельная библиотека; интерфейс (каталог \texttt{src/gui}) только вызывает
её методы.

\begin{center}
\begin{tabular}{p{5.2cm}p{9.8cm}}
\toprule
Файл & Содержимое \\
\midrule
\texttt{Time.h, Time.cpp} & модельное время, классы \texttt{Clock}, \texttt{WorkSchedule} \\
\texttt{Company.h, Company.cpp} & \texttt{Person}, \texttt{Employee}, \texttt{DepartmentHead},
  \texttt{Room}, \texttt{DepartmentRoom}, \texttt{ConferenceHall}, \texttt{Department},
  \texttt{Company} \\
\texttt{Events.h, Events.cpp} & \texttt{Event} и наследники: \texttt{DailyMeeting},
  \texttt{ManagersMeeting}, \texttt{Seminar}, \texttt{WorkMeeting}, \texttt{BusinessTrip} \\
\texttt{Calendar.h, Calendar.cpp} & \texttt{Calendar}, \texttt{EventFilter} \\
\texttt{Conflicts.h, Conflicts.cpp} & \texttt{Conflict}, \texttt{ConflictChecker},
  \texttt{FixOption}, \texttt{ConflictResolver} \\
\texttt{Notifier.h, Notifier.cpp} & \texttt{ReminderRule}, \texttt{Message}, \texttt{Notifier} \\
\texttt{Generator.h, Generator.cpp} & \texttt{EventGenerator} \\
\texttt{Simulation.h, Simulation.cpp} & \texttt{SimulationParams}, \texttt{Statistics},
  \texttt{PendingRequest}, \texttt{Simulation} \\
\midrule
\texttt{main.cpp} & главный файл, запуск приложения \\
\texttt{MainWindow.h, .cpp} & главное окно \texttt{MainWindow} \\
\texttt{ParamsDialog.h, .cpp} & окно параметров моделирования \\
\texttt{EventDialog.h, .cpp} & окно добавления и изменения события \\
\texttt{ConflictDialog.h, .cpp} & окно разрешения конфликта \\
\texttt{ReportDialog.h, .cpp} & итоговый отчёт \\
\texttt{CMakeLists.txt} & описание сборки \\
\bottomrule
\end{tabular}
\end{center}

% =====================================================================
\section{Пользовательский интерфейс}

При запуске открывается \textbf{окно параметров} (рис.~\ref{fig:params}), в котором
задаются параметры эксперимента: число отделов и сотрудников, длина периода, шаг,
среднее число заявок в день, случайное зерно, режим разрешения конфликтов и категория пользователя.
Кнопка <<Начать>> создаёт модель и открывает главное окно.

\begin{figure}[H]
\centering
\includegraphics[width=0.5\textwidth]{img/params.png}
\caption{Окно параметров моделирования}
\label{fig:params}
\end{figure}

\textbf{Главное окно} (рис.~\ref{fig:main}) содержит:
\begin{itemize}
  \item кнопки управления временем: <<Старт>> --- запуск (каждые 0,5~с выполняется
        один шаг), <<Шаг>> --- один шаг вручную, <<Стоп>> --- пауза; рядом
        показывается состояние (<<Идёт>>, <<Пауза>>, <<Завершено>>) и текущее
        модельное время;
  \item фильтр просмотра: дата, отдел, сотрудник; при движении времени дата
        переключается на текущий день;
  \item список событий выбранного дня и панель с подробностями выбранного события
        (время, место, важность, участники, число напоминаний);
  \item кнопки <<Добавить>>, <<Изменить>>, <<Удалить>> (две последние доступны
        только администратору);
  \item журнал: удалённые события, напоминания, заявки, конфликты;
  \item меню <<Эксперимент>>: новый эксперимент, итоговый отчёт, выход.
\end{itemize}

\begin{figure}[H]
\centering
\includegraphics[width=\textwidth]{img/main.png}
\caption{Главное окно}
\label{fig:main}
\end{figure}

\textbf{Окно события} (рис.~\ref{fig:event}) служит для добавления и изменения
события: вид, название, дата и время, длительность (для командировки --- число
дней), помещение, важность и участники. Проверяется, что выбраны участники,
событие лежит в рабочем времени и не в прошлом.

\textbf{Окно конфликта} (рис.~\ref{fig:conflict}) открывается, когда у заявки
обнаружен конфликт. В нём перечислены конфликты и предложенные системой варианты;
варианты, меняющие уже назначенные события, обычному пользователю недоступны.
Можно применить вариант, изменить событие вручную или отклонить заявку.

\begin{figure}[H]
\centering
\begin{minipage}{0.45\textwidth}
\centering
\includegraphics[width=\textwidth]{img/event.png}
\caption{Окно события}
\label{fig:event}
\end{minipage}\hfill
\begin{minipage}{0.52\textwidth}
\centering
\includegraphics[width=\textwidth]{img/conflict.png}
\caption{Окно конфликта}
\label{fig:conflict}
\end{minipage}
\end{figure}

По окончании периода (или из меню) открывается \textbf{итоговый отчёт}
(рис.~\ref{fig:report}): статистика и список состоявшихся событий по дням.
Отчёт можно сохранить в текстовый файл.

\begin{figure}[H]
\centering
\includegraphics[width=0.6\textwidth]{img/report.png}
\caption{Итоговый отчёт}
\label{fig:report}
\end{figure}

% =====================================================================
\section{Проведённые эксперименты}

Эксперименты проводились в режиме автоматического разрешения конфликтов, чтобы
моделирование шло без остановок. Каждый результат --- среднее по пяти запускам
с разными значениями параметра <<Случайное зерно>> (1--5).

\subsection*{Эксперимент 1. Интенсивность заявок}
$K = 6$, $N = 14$ дней, шаг 30 минут, администратор.

\begin{center}
\begin{tabular}{rrrrrr}
\toprule
$\lambda$ & заявок & с конфликтом & отклонено & загрузка зала & напоминаний \\
\midrule
1 & 7,8  & 2,4 (31\,\%)  & 0,0 & 11,7\,\% & 469 \\
3 & 22,6 & 7,4 (33\,\%)  & 0,0 & 23,7\,\% & 584 \\
5 & 37,8 & 20,4 (54\,\%) & 0,6 & 36,4\,\% & 780 \\
8 & 59,2 & 35,2 (59\,\%) & 2,0 & 46,9\,\% & 953 \\
\bottomrule
\end{tabular}
\end{center}

С ростом числа заявок доля заявок с конфликтами растёт с 31\,\% до 59\,\%:
чаще всего конфликт возникает из-за занятости конференц-зала, загрузка которого
выросла в 4 раза. Почти все конфликты удаётся разрешить, отклоняются единичные
заявки --- только при большой нагрузке, когда свободного времени не остаётся.

\subsection*{Эксперимент 2. Число отделов}
$\lambda = 5$, $N = 14$ дней.

\begin{center}
\begin{tabular}{rrrrr}
\toprule
$K$ & состоялось событий & с конфликтом & загрузка зала & напоминаний \\
\midrule
5 & 85,0  & 9,6 (33\,\%)  & 25,0\,\% & 525 \\
7 & 108,2 & 15,0 (46\,\%) & 30,7\,\% & 734 \\
9 & 131,4 & 14,6 (41\,\%) & 30,8\,\% & 891 \\
\bottomrule
\end{tabular}
\end{center}

С ростом числа отделов растёт число событий (больше ежедневных планёрок) и
напоминаний. Доля конфликтов растёт не монотонно: она зависит в основном от
занятости общего конференц-зала, а не от числа отделов.

\subsection*{Эксперимент 3. Категория пользователя}
$K = 6$, $N = 30$ дней.

\begin{center}
\begin{tabular}{llrrr}
\toprule
$\lambda$ & пользователь & с конфликтом & отклонено & перенесено назначенных \\
\midrule
8  & администратор & 100,6 & 1,2 & 1,6 \\
8  & обычный       & 99,8  & 2,2 & 0,0 \\
10 & администратор & 140,8 & 3,6 & 1,4 \\
10 & обычный       & 140,6 & 3,2 & 0,0 \\
\bottomrule
\end{tabular}
\end{center}

В автоматическом режиме выбирается первый допустимый вариант, обычно перенос
нового события, поэтому право администратора переносить уже назначенные события
используется редко (1--2 раза за месяц) --- только когда других вариантов нет.
Разница в числе отклонённых заявок невелика и сравнима со случайным разбросом.

\end{document}
